% Part: history
% Chapter: biographies
% Section: rozsa-peter

\documentclass[../../../include/open-logic-section]{subfiles}

\begin{document}

\olfileid{his}{bio}{pet}

\olsection{روژا پیتر}

\olphoto{peter-rozsa}{R\'ozsa P\'eter}

روژا پیتر با نام روژا پولیتسر، R\'ozsa Politzer، در 17 فوریهٔ 1905 در
بوداپست مجارستان به دنیا آمد. او بیش از همه برای کارش دربارهٔ تابع‌های
بازگشتی شناخته می‌شود؛ کاری که برای شکل‌گیری حوزهٔ نظریهٔ بازگشت اساسی
بود.

پیتر در دوره‌ای پرآشوب از نظر سیاسی بزرگ شد؛ هنگامی که نوجوان بود جنگ
جهانی اول در جریان بود. بااین‌حال توانست در مدرسهٔ مرفه دخترانهٔ ماریا
ترزیا در بوداپست تحصیل کند و در سال 1922 از آنجا فارغ‌التحصیل شد. سپس
در دانشگاه پازمانی پیتر در بوداپست، که بعدها دانشگاه اتووش لوراند
(E\"otv\"os Lor\'and University) نام گرفت، تحصیل کرد. به اصرار پدرش
تحصیل شیمی را آغاز کرد، اما بعد به ریاضیات تغییر رشته داد و در 1927
فارغ‌التحصیل شد. با آنکه صلاحیت تدریس ریاضی دبیرستان را داشت، وضعیت
اقتصادی آن زمان وخیم بود، زیرا رکود بزرگ بر اقتصاد جهان اثر گذاشته
بود. پیتر در این دوره کارهای موقت گوناگونی به‌عنوان آموزگار خصوصی و
معلم ریاضی انجام داد. سرانجام برای تحصیلات تکمیلی در ریاضیات به
دانشگاه بازگشت. ابتدا قصد داشت در نظریهٔ اعداد کار کند، اما وقتی فهمید
نتایجش پیش‌تر اثبات شده‌اند، نزدیک بود ریاضیات را یکسره کنار بگذارد.
او را به کار روی قضیه‌های ناتمامیت گودل تشویق کردند و بی‌آنکه بداند،
چند نتیجهٔ گودل را با روش‌هایی متفاوت اثبات کرد. این امر اعتمادبه‌نفسش
را بازگرداند و پیتر با الهام از برنامهٔ مبانی هیلبرت نخستین مقاله‌هایش
دربارهٔ نظریهٔ بازگشت را نوشت. در 1935 دکتری گرفت و در 1937 سردبیر
\emph{Journal of Symbolic Logic} شد.

مقاله‌های اولیهٔ پیتر به‌طور گسترده از آثار بنیان‌گذار نظریهٔ تابع‌های
بازگشتی دانسته می‌شوند. او در \cite{Peter1935a} رابطهٔ میان گونه‌های
مختلف بازگشت را بررسی کرد. در \cite{Peter1935b} نشان داد تابعی معین که
به‌طور بازگشتی تعریف شده است، بازگشتی اولیه نیست. این کار نتیجهٔ
پیشین ویلهلم آکرمان را ساده کرد. تابع ساده‌شدهٔ پیتر همان است که اکنون
اغلب تابع آکرمان، و گاهی دقیق‌تر تابع آکرمان--پیتر، نامیده می‌شود. او
نخستین کتاب دربارهٔ نظریهٔ تابع‌های بازگشتی را نوشت \citep{Peter1951}.

پیتر با وجود اهمیت و نفوذ کارش تا سال 1945 جایگاه تدریس تمام‌وقت به
دست نیاورد. در دوران اشغال مجارستان به دست نازی‌ها در جنگ جهانی دوم،
قوانین یهودستیزانه اجازهٔ تدریس به او نمی‌دادند. دولت در 1944 محله‌ای
یهودی‌نشین در بوداپست ایجاد کرد؛ این محله از بقیهٔ شهر جدا و با
نگهبانان مسلح محاصره شد. پیتر ناچار بود تا آزادسازی آن در 1945 در این
محله زندگی کند. سپس در کالج تربیت معلم بوداپست و از 1955 به بعد در
دانشگاه اتووش لوراند تدریس کرد. او نخستین زن ریاضی‌دان مجار بود که
عنوان دکتر آکادمی در ریاضیات را به دست آورد و نخستین زنی بود که به
عضویت آکادمی علوم مجارستان انتخاب شد.

پیتر به‌عنوان آموزگاری پرشور در ریاضیات شناخته می‌شد که ترجیح می‌داد
به جای صرفاً سخنرانی‌کردن، ماهیت و زیبایی مسئله‌های ریاضی را همراه با
دانشجویانش کشف کند. از همین رو دانشجویانش با محبت او را «خاله روژا»
می‌نامیدند. پیتر در سال 1977 و در 71سالگی درگذشت.

\begin{reading}
برای مطالب زندگی‌نامه‌ای بیشتر، \citep{Oconnor2014} و
\citep{Andrasfai1986} را ببینید. \citet{Tamassy1994} گفت‌وگویی کوتاه
با پیتر انجام داده است. برای نوشته‌ای سرگرم‌کننده دربارهٔ ریاضیات،
کتاب پیتر با عنوان \emph{Playing With Infinity} را ببینید
\citep{Peter2010}.
\end{reading}

\end{document}
